% GENERATED FILE. Edit canonical lessons, then run npm run generate:books.
% canonical-source-hash: fnv1a64:38a74deb6a1c7e23
% canonical-lessons: PT-C498-conviver, PT-C498-convivio, PT-C498-vizinhanca, PT-C498-confidente, PT-C498-lealdade, PT-R498-revisao-491-498

\chapter{Amigos e vizinhos}
\label{ch:pt-amigos-e-vizinhos}
\hlchaptermodality{498}

\begin{chapteropening}
\textbf{By the end of this chapter:} \emph{I can talk about getting on with friends and neighbours.}

Complete the last lesson of chapter 498: I can talk about getting on with friends and neighbours.
\end{chapteropening}

\section[conviver]{conviver — to socialise; to live together}

\label{lesson:PT-C498-conviver}



\begin{quote}
Before the new one: say the Portuguese for betrayal, then the Portuguese for to reconcile.
\end{quote}

\subsection*{You'll want to know: conviver}
\textbf{conviver} — \textquotedblleft{}to socialise; to live together\textquotedblright{}.

\begin{quote}
Gosto de conviver com os colegas depois do trabalho.
\end{quote}

\begin{quote}
\emph{I like socialising with my colleagues after work.}
\end{quote}

\emph{Conviver com} — to get on with, to live alongside.

\begin{cousinweb}
\emph{conviver} is a learned borrowing from Latin \emph{convīvere} \textquotedblleft{}to live together, to dine together\textquotedblright{} (Lewis \& Short), built on \emph{vīvere} \textquotedblleft{}to live\textquotedblright{}.
\end{cousinweb}

\begin{grammarlens}[title={What you've built}]
One more word for telling, explaining and describing at B1.
\end{grammarlens}

\subsection*{Your turn}
\begin{itemize}
\raggedright
  \item \emph{Say it:} \emph{conviver}
  \item \emph{Say it:} \emph{conviver}, once more
  \item \emph{Say it:} a sentence of your own with \emph{conviver}
  \item \emph{Recall:} say the Portuguese for betrayal, then the Portuguese for to reconcile, then say \emph{conviver} again
\end{itemize}

\begin{tcolorbox}[breakable,colback=teal!4,colframe=teal!35!black,title={Before you move on}]
What does conviver mean? (\textquotedblleft{}To socialise; to live together\textquotedblright{}.) Say it once more.
\end{tcolorbox}
\section[o convívio]{o convívio — get-together; social life}

\label{lesson:PT-C498-convivio}



\begin{quote}
Before the new one: say the Portuguese for to reconcile, then the Portuguese for to socialise.
\end{quote}

\subsection*{You'll want to know: o convívio}
\textbf{o convívio} — \textquotedblleft{}get-together; social life\textquotedblright{}.

\begin{quote}
No sábado há um convívio na associação do bairro.
\end{quote}

\begin{quote}
\emph{On Saturday there is a get-together at the residents' association.}
\end{quote}

In Portugal \emph{um convívio} is often a party or a shared meal.

\begin{cousinweb}
\emph{convívio} is a learned borrowing from Latin \emph{convīvium} \textquotedblleft{}a feast, a dinner party\textquotedblright{} (Lewis \& Short), built on \emph{vīvere} \textquotedblleft{}to live\textquotedblright{}.
\end{cousinweb}

\begin{grammarlens}[title={What you've built}]
One more word for telling, explaining and describing at B1.
\end{grammarlens}

\subsection*{Your turn}
\begin{itemize}
\raggedright
  \item \emph{Say it:} \emph{o convívio}
  \item \emph{Say it:} \emph{o convívio}, once more
  \item \emph{Say it:} a sentence of your own with \emph{o convívio}
  \item \emph{Recall:} say the Portuguese for to reconcile, then the Portuguese for to socialise, then say \emph{o convívio} again
\end{itemize}

\begin{tcolorbox}[breakable,colback=teal!4,colframe=teal!35!black,title={Before you move on}]
What does o convívio mean? (\textquotedblleft{}Get-together; social life\textquotedblright{}.) Say it once more.
\end{tcolorbox}
\section[a vizinhança]{a vizinhança — neighbourhood; the neighbours}

\label{lesson:PT-C498-vizinhanca}



\begin{quote}
Before the new one: say the Portuguese for to socialise, then the Portuguese for get-together.
\end{quote}

\subsection*{You'll want to know: a vizinhança}
\textbf{a vizinhança} — \textquotedblleft{}neighbourhood; the neighbours\textquotedblright{}.

\begin{quote}
Toda a vizinhança veio à festa.
\end{quote}

\begin{quote}
\emph{All the neighbours came to the party.}
\end{quote}

Built on \emph{o vizinho}.

\begin{cousinweb}
\emph{vizinhança} is built on \emph{vizinho} \textquotedblleft{}neighbour\textquotedblright{}, which is inherited from Latin \emph{vīcīnus} \textquotedblleft{}near, neighbouring\textquotedblright{} (Lewis \& Short).
\end{cousinweb}

\begin{grammarlens}[title={What you've built}]
One more word for telling, explaining and describing at B1.
\end{grammarlens}

\subsection*{Your turn}
\begin{itemize}
\raggedright
  \item \emph{Say it:} \emph{a vizinhança}
  \item \emph{Say it:} \emph{a vizinhança}, once more
  \item \emph{Say it:} a sentence of your own with \emph{a vizinhança}
  \item \emph{Recall:} say the Portuguese for to socialise, then the Portuguese for get-together, then say \emph{a vizinhança} again
\end{itemize}

\begin{tcolorbox}[breakable,colback=teal!4,colframe=teal!35!black,title={Before you move on}]
What does a vizinhança mean? (\textquotedblleft{}Neighbourhood; the neighbours\textquotedblright{}.) Say it once more.
\end{tcolorbox}
\section[o confidente]{o confidente — confidant}

\label{lesson:PT-C498-confidente}



\begin{quote}
Before the new one: say the Portuguese for get-together, then the Portuguese for neighbourhood.
\end{quote}

\subsection*{You'll want to know: o confidente}
\textbf{o confidente} — \textquotedblleft{}confidant\textquotedblright{}.

\begin{quote}
O meu melhor amigo é o meu confidente.
\end{quote}

\begin{quote}
\emph{My best friend is my confidant.}
\end{quote}

A woman is \emph{a confidente}.

\begin{cousinweb}
\emph{confidente} is a learned borrowing, perhaps through Italian, from Latin \emph{cōnfīdēns} \textquotedblleft{}trusting\textquotedblright{} (Lewis \& Short), from \emph{cōnfīdere} \textquotedblleft{}to trust fully\textquotedblright{}, built on \emph{fīdere} \textquotedblleft{}to trust\textquotedblright{}.
\end{cousinweb}

\begin{grammarlens}[title={What you've built}]
One more word for telling, explaining and describing at B1.
\end{grammarlens}

\subsection*{Your turn}
\begin{itemize}
\raggedright
  \item \emph{Say it:} \emph{o confidente}
  \item \emph{Say it:} \emph{o confidente}, once more
  \item \emph{Say it:} a sentence of your own with \emph{o confidente}
  \item \emph{Recall:} say the Portuguese for get-together, then the Portuguese for neighbourhood, then say \emph{o confidente} again
\end{itemize}

\begin{tcolorbox}[breakable,colback=teal!4,colframe=teal!35!black,title={Before you move on}]
What does o confidente mean? (\textquotedblleft{}Confidant\textquotedblright{}.) Say it once more.
\end{tcolorbox}
\section[a lealdade]{a lealdade — loyalty}

\label{lesson:PT-C498-lealdade}



\begin{quote}
Before the new one: say the Portuguese for neighbourhood, then the Portuguese for confidant.
\end{quote}

\subsection*{You'll want to know: a lealdade}
\textbf{a lealdade} — \textquotedblleft{}loyalty\textquotedblright{}.

\begin{quote}
Admiro a lealdade do meu cão.
\end{quote}

\begin{quote}
\emph{I admire my dog's loyalty.}
\end{quote}

The adjective is \emph{leal}.

\begin{cousinweb}
\emph{lealdade} is inherited, through Old Portuguese, from Medieval Latin \emph{lēgālitās} (EtymDB), built on \emph{lēgālis} \textquotedblleft{}of the law\textquotedblright{}, from \emph{lēx} \textquotedblleft{}law\textquotedblright{}.
\end{cousinweb}

\begin{grammarlens}[title={What you've built}]
One more word for telling, explaining and describing at B1.
\end{grammarlens}

\subsection*{Your turn}
\begin{itemize}
\raggedright
  \item \emph{Say it:} \emph{a lealdade}
  \item \emph{Say it:} \emph{a lealdade}, once more
  \item \emph{Say it:} a sentence of your own with \emph{a lealdade}
  \item \emph{Recall:} say the Portuguese for neighbourhood, then the Portuguese for confidant, then say \emph{a lealdade} again
\end{itemize}

\begin{tcolorbox}[breakable,colback=teal!4,colframe=teal!35!black,title={Before you move on}]
What does a lealdade mean? (\textquotedblleft{}Loyalty\textquotedblright{}.) Say it once more.
\end{tcolorbox}
\section[(dialogue)]{Revisão dos capítulos 491–498}

\label{lesson:PT-R498-revisao-491-498}



\begin{quote}
Say the Portuguese for to date and relationship.
\end{quote}

\subsection*{Your turn}
Say these aloud:

\begin{itemize}
\raggedright
  \item to date, relationship, passion, to fall in love, to hug
  \item brother-in-law, stepfather, stepmother, stepson, relative
  \item great-grandfather, great-grandson, orphan, to adopt, adoption
  \item nappy, baby's bottle, dummy, to crawl, to rock
  \item to bring up, to behave, to punish, (monthly) pocket money, spoiled
  \item tenderness, affectionate, affection, to stroke, to spoil
  \item to quarrel, to separate, to betray, betrayal, to reconcile
  \item to socialise, get-together, neighbourhood, confidant, loyalty
\end{itemize}

\begin{tcolorbox}[breakable,colback=teal!4,colframe=teal!35!black,title={Before you move on}]
To date? (\textbf{namorar}.) Relationship? (\textbf{o namoro}.) Passion? (\textbf{a paixão}.) To fall in love? (\textbf{apaixonar}.) To hug? (\textbf{abraçar}.) Brother-in-law? (\textbf{o cunhado}.) Stepfather? (\textbf{o padrasto}.) Stepmother? (\textbf{a madrasta}.) Stepson? (\textbf{o enteado}.) Relative? (\textbf{o parente}.) Great-grandfather? (\textbf{o bisavô}.) Great-grandson? (\textbf{o bisneto}.) Orphan? (\textbf{o órfão}.) To adopt? (\textbf{adotar}.) Adoption? (\textbf{a adoção}.) Nappy? (\textbf{a fralda}.) Baby's bottle? (\textbf{o biberão}.) Dummy? (\textbf{a chupeta}.) To crawl? (\textbf{gatinhar}.) To rock? (\textbf{embalar}.) To bring up? (\textbf{educar}.) To behave? (\textbf{comportar}.) To punish? (\textbf{castigar}.) (monthly) pocket money? (\textbf{a mesada}.) Spoiled? (\textbf{mimado}.) Tenderness? (\textbf{a ternura}.) Affectionate? (\textbf{carinhoso}.) Affection? (\textbf{o afeto}.) To stroke? (\textbf{acariciar}.) To spoil? (\textbf{mimar}.) To quarrel? (\textbf{brigar}.) To separate? (\textbf{separar}.) To betray? (\textbf{trair}.) Betrayal? (\textbf{a traição}.) To reconcile? (\textbf{reconciliar}.) To socialise? (\textbf{conviver}.) Get-together? (\textbf{o convívio}.) Neighbourhood? (\textbf{a vizinhança}.) Confidant? (\textbf{o confidente}.) Loyalty? (\textbf{a lealdade}.)
\end{tcolorbox}
